\chapter{Antibiotic Resistance}
\index{Antibiotic Resistance}

\section*{Definition and Overview}
Antibiotic resistance occurs when bacteria change in ways that make the antibiotics designed to kill them ineffective. Resistant bacteria survive exposure to a drug and continue to multiply, so infections that were once easily treated become prolonged, more severe, or occasionally untreatable. The term \textbf{antimicrobial resistance} (AMR) is used more broadly to describe the same phenomenon in viruses, fungi, and parasites, and the World Health Organization ranks AMR among the top ten global public health threats facing humanity.

Resistance is a natural evolutionary process, but human activity -- particularly the overuse and misuse of antibiotics in human medicine, agriculture, and animal husbandry -- has accelerated it dramatically. The consequences reach far beyond the treatment of individual infections: routine operations, cancer chemotherapy, organ transplantation, and care of premature babies all depend on effective antibiotics to prevent and treat infection. When antibiotics fail, these procedures become far riskier.

\section*{Epidemiology}
Resistant infections are now documented in every country, although rates vary widely. Several common organisms have already become resistant to multiple drugs, including methicillin-resistant \textit{Staphylococcus aureus} (MRSA), extended-spectrum beta-lactamase (ESBL)-producing \textit{Escherichia coli}, multidrug-resistant tuberculosis, carbapenem-resistant \textit{Pseudomonas}, and strains of gonorrhoea, typhoid, and pneumococcus. Infections caused by resistant bacteria are estimated to have contributed to around 1.27 million deaths worldwide in 2019 alone, with many more deaths in which resistance played a contributing role.

A major driver is prescribing: a large proportion of antibiotics given in the community are for conditions that would resolve on their own, particularly colds and other viral respiratory infections, which do not respond to antibiotics at all. Broad-spectrum antibiotics are frequently used when a narrow-spectrum drug would suffice, and courses are often continued for longer than necessary. In low- and middle-income countries, limited laboratory capacity, counterfeit medicines, and restricted access to second-line drugs compound the problem.

\section*{Aetiology and Pathophysiology}
Bacteria acquire resistance in two principal ways: through spontaneous \textbf{mutation} of their own chromosomal genes, and through \textbf{horizontal gene transfer}, in which bacteria share pieces of DNA (plasmids, transposons, or integrons) carrying resistance genes with unrelated bacteria, even across species. Mobile genetic elements can carry resistance to several classes of antibiotic at once, which is why multidrug resistance spreads so quickly.

When an antibiotic is used, it kills susceptible bacteria but leaves resistant ones behind, which then multiply and become dominant -- a process called \textbf{selection pressure}. The mechanisms by which bacteria defend themselves include:

\begin{itemize}
    \item \textbf{Enzymatic inactivation:} bacteria produce enzymes, such as beta-lactamases, that destroy or chemically modify the antibiotic before it acts
    \item \textbf{Efflux pumps:} transport proteins actively pump the antibiotic out of the bacterial cell, keeping intracellular concentrations too low to be effective
    \item \textbf{Altered drug targets:} bacteria change the molecular structure of the protein or ribosome the antibiotic normally binds to, so the drug can no longer attach
    \item \textbf{Reduced permeability:} bacteria thicken or alter their outer membrane so that the antibiotic cannot penetrate the cell
    \item \textbf{Alternative pathways:} bacteria bypass the metabolic step the antibiotic blocks and use a different route to the same product
    \item \textbf{Biofilms:} bacteria living in dense surface films, as on catheters, heart valves, or wounds, are shielded from antibiotics and host defences
\end{itemize}

\section*{Clinical Features}
Resistant infections do not produce unique symptoms; they present with the features of the underlying infection, but these tend to be more severe, longer lasting, or slower to resolve. Typical patterns include:

\begin{itemize}
    \item An infection that fails to improve, or worsens, after several days of an appropriate antibiotic
    \item Recurrent infections that require repeated or escalating courses of treatment
    \item Persistent fever, pain, purulent discharge, or malaise beyond the expected course of illness
    \item Surgical site, urinary, or bloodstream infections in hospitalised patients exposed to broad-spectrum antibiotics
    \item The carriage of resistant bacteria without any symptoms at all -- colonisation of the skin, gut, or nose that can spread to others
\end{itemize}

\section*{Differential Diagnosis}
When an infection is not responding as expected, several possibilities other than antibiotic resistance must be considered:

\begin{itemize}
    \item The \textbf{wrong diagnosis} -- the illness may be viral, fungal, inflammatory, or non-infectious rather than bacterial
    \item An \textbf{inappropriate antibiotic} for the infecting organism, or a dose that is too low
    \item \textbf{Poor adherence}, missed doses, or \textbf{poor absorption} of the medicine
    \item A \textbf{new or superimposed infection} with a different organism
    \item An \textbf{undrained focus} of infection, such as an abscess, foreign body, or obstructed tube, which requires physical source control rather than more antibiotics
    \item \textbf{Non-infectious causes} of fever, such as drug reactions, thrombosis, or inflammatory disease
\end{itemize}

\section*{When to Seek Medical Care}
Return to a doctor if an infection is not improving after a few days of antibiotics, if symptoms worsen, or if fever persists. Seek review early for infections in young infants, the elderly, pregnant women, and people with weakened immune systems or chronic disease.

\textbf{Contact your local emergency services for:}
\begin{itemize}
    \item Difficulty breathing, chest pain, or a bluish tinge to the lips or skin
    \item Confusion, drowsiness, or difficulty rousing
    \item A high fever with a non-blanching rash, stiff neck, or sensitivity to light
    \item Signs of sepsis: rapid breathing, rapid heart rate, mottled or cold skin, or greatly reduced urine output
    \item Severe pain, spreading redness, or a wound that is discharging pus or blood
\end{itemize}

\section*{Diagnosis and Investigations}
The assessment of a suspected resistant infection begins with a careful history of the illness, prior antibiotic courses, recent hospitalisation or travel, and exposure to others with resistant infections. Examination focuses on the site of infection, vital signs, and signs of systemic illness or sepsis. Laboratory and radiological investigations include:

\begin{itemize}
    \item \textbf{Microscopy and culture:} specimens such as blood, urine, sputum, or wound swabs are grown to identify the organism
    \item \textbf{Susceptibility (sensitivity) testing:} the bacterium is exposed to a panel of antibiotics to determine which drugs are effective
    \item \textbf{Molecular (PCR) testing:} rapid tests that detect specific bacteria and resistance genes, such as \textit{mecA} (MRSA) or \textit{blaCTX-M} (ESBL), within hours
    \item \textbf{Antibiograms:} local summaries showing which antibiotics are most likely to work for common organisms, guiding empirical therapy
    \item \textbf{Inflammatory markers:} blood counts, C-reactive protein, and procalcitonin help gauge the severity and nature of infection
    \item \textbf{Imaging:} chest X-ray, ultrasound, or CT to identify the site and extent of infection
\end{itemize}

Treatment is usually started \textbf{empirically} -- based on the most likely organism and local resistance patterns -- and then refined, or \textbf{de-escalated}, once culture results are known.

\section*{Management}
The management of resistant infection rests on both treating the individual patient and protecting the antibiotic supply for the community:

\begin{itemize}
    \item \textbf{Antibiotic stewardship:} prescribing the narrowest effective drug, at the right dose and route, for the shortest adequate course, and stopping antibiotics when infection is excluded
    \item \textbf{Sensitivity-directed therapy:} switching to the drug shown by culture to be active against the organism
    \item \textbf{Combination therapy and newer agents:} for serious resistant infections, using two drugs together or reserving last-line agents such as carbapenems, tigecycline, or newer cephalosporins with beta-lactamase inhibitors
    \item \textbf{Source control:} draining abscesses, removing infected catheters, or debriding necrotic tissue, since antibiotics alone cannot cure an undrained collection
    \item \textbf{Infection prevention:} isolation of patients with resistant organisms, hand hygiene, and contact precautions in hospitals
    \item \textbf{Hospital care:} intravenous therapy, organ support, and monitoring for severe or invasive infections
\end{itemize}

\section*{Complications}
\begin{itemize}
    \item Prolonged illness, treatment failure, and increased rates of serious complications such as sepsis and death
    \item Greater reliance on stronger, more toxic, or more expensive last-line drugs with additional side effects
    \item Longer hospital stays, more surgical procedures, and higher healthcare costs
    \item Transmission of resistant bacteria within households, hospitals, and the wider community
    \item Loss of effective prophylaxis, undermining routine surgery, cancer chemotherapy, and transplant medicine
    \item The emergence of pan-resistant organisms for which no reliable treatment exists
\end{itemize}

\section*{Prognosis and Outcomes}
Most resistant infections are still treatable, but recovery typically takes longer, may require intravenous or combination therapy, and carries a higher risk of complications. Outcomes depend on the organism, the site and severity of infection, and the patient's underlying health and immune status. For infections caused by pan-resistant organisms, mortality is substantial.

On a population level, resistance can be slowed and in some settings reversed through disciplined prescribing, infection prevention, vaccination, and the development of new antibiotics. National action plans coordinated by the World Health Organization aim to reduce the global burden, but the effort depends on the prescribing decisions of clinicians and the expectations of patients alike. Protecting the antibiotics we currently have is a shared responsibility.

\section*{Prevention and Self-Care}
\begin{itemize}
    \item Use antibiotics only when a doctor confirms they are needed for a bacterial infection
    \item Do not demand antibiotics for colds, flu, or other viral illnesses
    \item Never share antibiotics with others or keep leftover tablets for a future illness
    \item Take antibiotics exactly as prescribed for the prescribed duration; do not share, save, extend, or restart them without medical advice
    \item Practise good hand hygiene and cover coughs and sneezes to reduce the spread of infection
    \item Keep vaccinations up to date, including flu and pneumococcal vaccines where recommended
    \item Prepare and store food safely, and handle animals and animal waste hygienically
    \item Support efforts by clinicians, farmers, and industry to reduce unnecessary antibiotic use
\end{itemize}

\section*{Sources and Further Reading}
The following organisations provide reliable, up-to-date information on antibiotic resistance for patients, students, and health professionals.

\begin{itemize}
    \item World Health Organization. \textit{Antimicrobial resistance.} https://www.who.int/health-topics/antimicrobial-resistance
    \item Centers for Disease Control and Prevention. \textit{Antibiotic resistance threats and stewardship resources.} https://www.cdc.gov/
    \item NHS. \textit{Information on antibiotics, resistance, and safe use.} https://www.nhs.uk/conditions/
    \item NICE. \textit{Guidance on antimicrobial stewardship and prescribing.} https://www.nice.org.uk/
    \item MedlinePlus. \textit{Health information on antibiotics and drug resistance.} https://medlineplus.gov/
    \item MSD Manual. \textit{Professional overview of antimicrobial resistance.} https://www.msdmanuals.com/
\end{itemize}
